\documentclass[11pt,a4paper]{article}
% Shared preamble for RuPhO English problem statements (match T1 2026 style).
% Use: \input{latex/rupho_en_preamble.tex} after \documentclass.
\usepackage[margin=1in]{geometry}
\usepackage{amsmath,amssymb}
\usepackage{graphicx}
\usepackage{float}
\usepackage{enumitem}
\usepackage{microtype}
\usepackage{caption}
\usepackage{tabularx}
\usepackage{hyperref}

\captionsetup{font=small,labelfont=bf,skip=6pt}

\setlist[itemize]{leftmargin=1.2cm}
\setlist[enumerate]{leftmargin=1.2cm}

% One outer box per subproblem: [id | statement | points] (no inner rules)
\newcommand{\problembox}[3]{%
  \par\addvspace{0.42em}%
  \noindent
  \setlength{\fboxsep}{7.5pt}%
  \setlength{\fboxrule}{0.95pt}%
  \fbox{%
    \begin{minipage}{\dimexpr\linewidth-2\fboxsep-2\fboxrule\relax}
    \begin{tabularx}{\linewidth}{@{}>{\raggedright\arraybackslash\bfseries}p{2.1em} @{\hspace{0.9em}} >{\raggedright\arraybackslash}X @{\hspace{0.9em}} >{\raggedleft\arraybackslash\bfseries}p{2.4em}@{}}
    #1 & #2 & #3
    \end{tabularx}
    \end{minipage}%
  }%
  \par\addvspace{0.32em}%
}

\title{\textbf{Analemma}}
\author{\normalsize Y18 (2018) --- English translation}
\date{}
\begin{document}
\maketitle

People of antiquity and the Middle Ages determined time by the position of the Sun in the sky and associated the timing of religious holidays and daily routines with it. Many ceremonies were performed at noon, which was defined as the time when the sun occupies its highest position in the sky. At this time, a ray of sun passing through a hole in the roof of a temple or cathedral crossed the meridian line on the floor, indicating the north-south direction (see figure). In English, the Latin designations for time before (a.m., ante meridiem) and after noon (p.m., post meridiem), originating at this time, have been preserved. By the angular altitude of the sun at noon (as well as by the position of the stars at midnight), travelers of the past, using special astronomical instruments, determined the geographic latitude of the ship.

\begin{figure}[H]
    \centering
    \includegraphics[width=0.72\textwidth]{T5_EN_assets/fig1.jpg}
    \caption{Figure 1.}
\end{figure}

The time reference system associated with the sun is called “true solar time” (TST) - in it the position of the sun is determined by the azimuth corresponding to the position of the sun on the celestial sphere. With the advent of accurate mechanical watches, it turned out to be easy to verify that true solar time is uneven - throughout the year, the length of the true solar day (the time between solar noons) changes. This is how the system of average solar time (AST) appeared - the length of the day in such a system coincides with the average length of the day throughout the year, time is counted uniformly.

\begin{figure}[H]
    \centering
    \includegraphics[width=0.72\textwidth]{T5_EN_assets/fig2.jpg}
    \caption{Figure 2.}
\end{figure}

If at a certain moment you synchronize the average solar time with the true one (at the moment of solar culmination, set the mechanical clock to noon), then during the year there will be a discrepancy between the time, determined in two different ways: for example, the sun at 12:00 local time on different days is either west or east of the meridional line. If you mark the position of the sun's ray at 12:00 during the year, instead of a straight line you will get a “figure eight” (see Fig. 1b). Obviously, this figure eight is associated with the position of the sun on the celestial sphere, and if, using more modern technologies, photographs of the sky are taken in a certain place in 12-hour increments, then the sun throughout the year also describes a curve shaped like a figure eight (see figure), which we will call an analemma.

\begin{figure}[H]
    \centering
    \includegraphics[width=0.72\textwidth]{T5_EN_assets/fig3.jpg}
    \caption{Figure 3.}
\end{figure}

The problem proposes to analyze various factors influencing the discrepancy between TST and MST, estimate its magnitude, and explain the shape of the analemma. Although the exact calculation of the desired discrepancy is complex and leads to cumbersome results, in practice this discrepancy is described with good accuracy by the sum of several sinusoids, the parameters of which we suggest you find. For simplicity, we will consider not the Earth but another very similar planet (see Fig.~1), rotating around its axis with a period of $\tau = 24~\text{h}$, with a year lasting $\tau N_0$. The axis of the planet's own rotation forms a small angle $\theta$ ($\theta \ll 1$) with the normal to the orbital plane. Consider $\theta$ given. Also assume that the eccentricity $e$ is given.

\textit{Fig.~1.} Position at $N=0$ on the first day of the year. The orbit (with eccentricity $e$) and rotation axis (inclined at an angle $\theta$ to the normal vector to the trajectory) of the hypothetical planet considered in the problem. The directions of rotation of the planet around the star and around its own axis coincide (counterclockwise when viewed from the north pole of the ecliptic).

\begin{figure}[H]
    \centering
    \includegraphics[width=0.72\textwidth]{T5_EN_assets/fig4.jpg}
    \caption{Figure 4.}
\end{figure}

We will assume that on the first day of the year ($N=0$) the planet is closest to the star, and the rotation axis lies in the plane specified by the planet-star radius vector and the normal to the trajectory plane. (For the Earth, this condition is not true, although the deviation is small.) The directions of rotation of the planet around the star and around its axis coincide (counterclockwise, when viewed from the “north” pole of the ecliptic). In the northern hemisphere, at the beginning of the year, days are shorter than nights.

Fig.~2. Angular system of celestial coordinates: $\varepsilon$ - height of the sun, $\Delta \lambda$ - azimuthal deviation from the average (south) direction to the midday sun.

\begin{figure}[H]
    \centering
    \includegraphics[width=0.72\textwidth]{T5_EN_assets/fig5.jpg}
    \caption{Figure 5.}
\end{figure}

Fig.~3. Longitude $\lambda$ and latitude $\varphi$ of a certain point on the surface of the planet.

\begin{figure}[H]
    \centering
    \includegraphics[width=0.72\textwidth]{T5_EN_assets/fig6.jpg}
    \caption{Figure 6.}
\end{figure}

\section*{Part A. Maximum altitude of the Sun}

\problembox{A1}{The greatest angular altitude of the sun during the day depends on the time of year and on the latitude at which the observer is located. Solve a problem worthy of an ancient navigator and find the coefficients $A_0$, $k_0$ and $\phi_0$ in an approximate expression describing the dependence of the deviation of this angular height (Fig. 2) from the value averaged over all days of the year ($N$ – day number). Consider the orbit to be circular. $$\Delta{\varepsilon}=A_0\sin\left(k_0N+\phi_0\right) $$}{0.90}

\problembox{A2}{Draw what curve a ray of light would describe over a year on the floor of the temple, the positions of which are measured every day at 12:00 TST, mark the position of the cardinal points in the drawing. (The cardinal directions on the planet in question are determined in the same way as on earth, in accordance with the direction of rotation of the planet around its axis).}{0.10}

\section*{Part B. Eccentric orbit with perpendicular axis}

In Part B, consider that the planet's axis is perpendicular to the orbital plane ($\theta=0$) and the orbit is close to a circle ($0<e\ll{1}$).

\problembox{B1}{Indicate at what point in the orbit the true solar day will be longest? In short?}{0.10}

\problembox{B2}{How long is the longest true day longer than the shortest?}{2.00}

\problembox{B3}{Assuming that on the first day of the year ($N=0$) the planet is closest to the star, find the parameters $A_1$, $k_1$ and $\phi_1$ of the approximate equation describing the difference in the duration of the true and average solar days depending on the number of the day in the year $N$ ($0\leq{N}\leq{364}$): $$t_\text{day TST} - t_\text{day MST} = A_1 \sin(k_1 N + \phi_1).$$}{0.60}

\problembox{B4}{What difference $\Delta_1(N)$ between TST and MST will accumulate over $N$ days if the systems were synchronized on the first day of the year?}{0.20}

\problembox{B5}{What is the largest discrepancy between the TST and the MST ($\Delta_{1max}$) during the year?}{0.10}

\section*{Part C. Tilt Axis in Circular Orbit}

Let now the planet move in a circular orbit ($e=0$), and the axis of its own rotation is inclined from the normal to the trajectory plane at a small angle $\theta$ ($\theta\ll{1}$). We will assume that on the first day of the year ($N=0$) the axis of proper rotation lies in the plane formed by the normal to the trajectory plane and the radius vector connecting the centers of the star and the planet. We will assume that the directions of rotation of the planet around the star and around its axis coincide (both movements are counterclockwise when viewed from the “north” pole of the ecliptic). In the northern hemisphere, at the beginning of the year, days are shorter than nights.

\problembox{C1}{Let us mark the point on the equator of the planet at which the star is at its zenith at some point in time. Let us assign zero longitude $\lambda=0$ to this point and monitor the points for which the star will be at the zenith after exactly 1 day of the MST, 2 days of the MST, etc. These points form a certain curve on the surface of the planet, which can be specified in terms of latitude $\varphi$ and longitude $\lambda$ (see Fig. 3). Get the formula for this curve.}{1.00}

\problembox{C2}{Suppose that at some moment when constructing the curve described in paragraph $\mathrm C1$, at a point with longitude $\lambda$ the star was at the zenith. Find $\Delta{\lambda}/\Delta{t}$, where $\Delta{\lambda}$ is the change in longitude of the point over time $\Delta{t}$ equal to one day.}{0.20}

\problembox{C3}{Find the parameters $A_2$, $k_2$ and $\phi_2$ of the approximate equation that best describe the difference between the duration of the true and average solar days on such a planet: $$t_\text{day TST} - t_\text{day MST} = A_2 \sin(k_2 N + \phi_2).$$}{1.50}

\problembox{C4}{What difference $\Delta_2(N)$ between the TST and the MST will accumulate over $N$ days, if on the first day of the year both systems were synchronized?}{0.20}

\problembox{C5}{What value $\Delta_{2max}$ does the greatest discrepancy between the TST and MST reach during the year?}{0.10}

\section*{Part D. Combination of contributions}

In Part D, for simplicity, we will assume that the contributions from the inclination of the axis and from the ellipticity of the orbit to the angular coordinates of the sun on the celestial sphere (see figure) add up.

\problembox{D1}{Obtain the dependences of both angular coordinates of the sun $\varepsilon(N)$ and $\Delta{\lambda}(N)$ at noon according to the MST (12:00) on the number of the day in the year $N$ at latitude $\varphi$.}{1.00}

\problembox{D2}{Draw schematically in the angular coordinates of the celestial sphere $\varepsilon$ and $\Delta{\lambda}$ (up to a constant) the analemmas for the cases described in parts B and C.}{1.00}

\problembox{D3}{Using pictures and formulas, describe a method that allows you to determine $e$ and $\theta$ from the parameters of the analemma.}{0.20}

Figure 4 below shows the analemmas of the four planets. The horizontal and vertical axes are $\Delta{\lambda}$ and $\varepsilon$ in degrees, respectively.

\begin{figure}[H]
    \centering
    \includegraphics[width=0.72\textwidth]{T5_EN_assets/fig7.jpg}
    \caption{Figure 7.}
\end{figure}

\problembox{D4}{Determine the numerical values of $e$ and $\theta$ for the planets whose analemmas are shown in the figure, and also indicate at what latitudes $\varphi$ the given curves were measured (all measurements were carried out in the northern hemisphere). The horizontal and vertical axes are $\Delta \lambda$ and $\varepsilon$ in degrees, respectively.}{0.80}

\vspace{1em}
\noindent\textit{Source:} \url{https://pho.rs/p/3045}

\end{document}